For this view, we are now standing on the Moon. At this instant viewed, the
Earth is centrally eclipsing the Sun (see the red circle on the sky).
Consequently the Moon is at one of its nodes now. Assume the longitudinal
and latitudinal librations are exactly $0^\circ$ at this moment.

\begin{description}
\item[(3.1)] At the time of this observation, which season is it in Hungary?
  (Circle the correct answer.) Spring / Summer / Autumn / Winter
\item[(3.2)] There is a yellow circle on the projected sky (next to the red
  circle), which denotes minor planet Juno, which is at a distance of
  exactly 3 au from the Sun at this moment. Estimate its distance to the
  Moon at this instant. (Rounded to the nearest integer in units of million
  km.) Assume all orbits to be circular.
\item[(3.3)] Approximately how much time (in Earth days) after the projected
  event will \ldots
  \begin{description}
  \item[\rm\ldots] the Sun set at your observing site?
  \item[\rm\ldots] the Earth set at your observing site?
  \end{description}
\item[(3.4)] Determine the Lunar (Selenographic) coordinates of this
  observing site (as defined in the provided lunar map). What is the name of
  the large surface lunar area where your observing site is situated? Do not
  use your national language, please use the official IAU nomenclature. (See
  the lunar map.)
\item[(3.5)] Estimate the distance from the observing site to the Apollo-11
  landing site ($0.6875^\circ$\,N, $23.4333^\circ$\,E).
\end{description}

% FIGURE NEEDED: lunar map defining the selenographic coordinate system,
% used for tasks 3.4 and 3.5
